\section{Conclusion}

This work presented a modular deep reinforcement learning framework 
for multi-asset financial trading, integrating custom trading environments, 
temporal observation processing, and multiple policy architectures, 
including MLP, Transformer-based, hybrid MLP--Transformer, and 
recurrent (LSTM) models.

The training analysis revealed clear differences in learning dynamics across 
algorithms. On-policy methods such as A2C showed limited learning capacity, 
with near-flat reward trajectories. PPO variants exhibited the most stable 
and consistent behavior across architectures, with Transformer-based policies 
achieving the strongest performance due to improved temporal modeling.

Off-policy methods were more sensitive to architectural design. DDPG (MLP) 
achieved the most stable and competitive performance overall. In contrast, 
Transformer-based variants of DDPG, TD3, and SAC showed instability and 
early termination, highlighting difficulties in combining replay-based 
learning with high-capacity sequence encoders.

The evaluation of financial metrics (total return and Sharpe ratio) was 
largely consistent with episodic reward trends, but revealed additional 
instability in long-horizon compounding behavior. In particular, 
recurrent PPO and Transformer-based TD3 exhibited episodic return spikes, 
indicating numerical instability and sensitivity to compounding effects in 
certain regimes.

The backtesting results further highlighted generalization gaps across 
architectures. While Transformer-based PPO models performed well during 
training, they failed to generalize to unseen market conditions, producing 
negative returns over the evaluation period. In contrast, DDPG (MLP), which
also achieved the strongest training stability, demonstrated the best 
out-of-sample performance with consistently positive returns.

To address these performance gaps, the implementation of a Mamdani-style fuzzy 
logic layer successfully replaced rigid, threshold-based constraints with smooth 
action adjustments and context-dependent risk penalties. Concurrently, the 
integration of a Genetic Algorithm provided an efficient evolutionary search 
mechanism to systematically discover robust hyperparameter configurations within 
highly non-linear training dynamics.

Overall, the results indicate a clear trade-off between expressiveness and 
stability. MLP-based policies provide the most robust and reliable performance, 
Transformer architectures improve representation power but reduce stability 
in off-policy settings, and recurrent architectures introduce additional 
instability without consistent performance gains in this experimental setup. 
However, by augmenting the core pipeline with layered fuzzy risk controls and genetic 
optimization, the framework establishes a stabilized training track that effectively 
safeguards portfolio capital without sacrificing architectural expressiveness.